\chapter{Introduction}

\section{Background of the Study}

When evaluating a company, investors, lenders, and regulators mostly rely on financial statements. Managers have some control over what they report under accrual accounting since reported profit is based on estimates and assessments regarding revenue recognition, receivables, inventories, and depreciation. Earnings management is the procedure of managers using their autonomy to alter reported earnings, either to deceive certain stakeholders about the company's underlying performance or to affect outcomes that rely on reported numbers, like bonuses or loan covenants \parencite{healy1999}.

Between managers and account readers are two defenses. The board's audit committee is in charge of financial reporting and the interaction with the external auditor, who assesses whether the financial statements adhere to the relevant reporting framework. Since it is impossible to directly witness how well these safeguards function, research examines characteristics that can be quantified, such as the length of time the same audit firm has worked with the client (audit tenure), the amount the customer pays for the audit (audit fees), and the independence of the audit committee. There is inconsistent research regarding the impact of each element, which can either increase or decrease the safety.

In practice, earnings management is important since its expenses are quantifiable. \textcite{ilyas2019} discover that companies that manipulate earnings have a higher cost of capital in a sample of 144 non-financial PSX companies from 2006 to 2016. This is because manipulated information makes investors less confident in reported fundamentals.

\textcite{bhutta2022} highlight inadequate regulatory implementation in the corporate sector and discover that real earnings management promotes investment inefficiency in Pakistani non-financial enterprises. The regulator has also taken measures over the protections themselves: cases under Regulation 27 on audit committee composition and Regulation 32 on auditor appointment are included in SECP enforcement actions issued between 2023 and 2025 \parencite{secp2023}.

This study examines whether audit tenure, audit fees, and audit committee independence are associated with the level of accrual-based earnings management in textile composite companies listed on the Pakistan Stock Exchange (PSX).

\section{Overview of the Textile Composite Sector and the Regulatory Setting}

Textile composite firms include many textile production steps, namely spinning, weaving, processing, and stitching, into a single enterprise. As a result, accruals are comparable between businesses since their financial statements include equivalent assets, including revenue, trade receivables, inventory, and property, plant, and equipment. The industry is vital to Pakistan's export-oriented economy. According to data from the Pakistan Bureau of Statistics published by \textcite{profit2026}, textile exports amounted to USD 17.93 billion in FY2025--2026, or roughly 59\% of overall exports of USD 30.14 billion. On October 25, 2026, \textcite{tickeranalysts2026} listed 37 businesses on PSX, while \textcite{finhisaab2026} listed 53. Market-data providers cannot agree on the size of the sector. As a result, the 37-company list serves as the study's population, which is detailed in Section~\ref{sec:sampling}.

Under the Companies Act of 2017, listed firms publish audited yearly financial statements, and the auditors' compensation is disclosed in the statements' notes. Rules for audit committees and external auditors are defined in the Listed Companies (Code of Corporate Governance) Regulations, 2019. They were notified by SECP on September 25, 2019, as SRO 1163(I)/2019, and they became operative upon publication \parencite{secp2019}. As per Regulation 27, the board must form an audit committee with a minimum of three non-executive directors, one independent director, and an independent chair who is not the board chair. The external auditors' appointment, dismissal, and audit fees are recommended by the audit committee to the board. Regulation 32 prohibits a business from selecting an audit firm that is either not registered with the Audit Oversight Board or has an unsatisfactory grade under the Institute of Chartered Accountants of Pakistan's Quality Control Review program. Regulation 33 mandates that firms in the financial industry replace their external auditors every five years. However, other listed companies are merely required to rotate their engagement partner every five years, unless the audit company is a sole proprietorship. Thus, a textile composite manufacturer can maintain the same audit firm for many years, which gives audit tenure real variation in this sector. The primary source of information for audit committee independence is the statement of compliance, which is required by Annexure A of the Regulations and includes a list of board members by category as well as the names of the audit committee members and chair.

\section{Problem Statement}

Earnings management undermines the credibility of reported profit, and investors and lenders rely on auditors and audit committees to limit it. The cost is substantial, as earnings manipulation enhances the cost of financing for Pakistani listed companies \parencite{ilyas2019}. It is uncertain whether longer audit tenure, more audit fees, and a more independent audit committee are associated with lower earnings management in Pakistani listed businesses. Other countries' findings are equivocal, and Pakistani studies on audit committee independence show conflicting results \parencite{latif2015,mumtaz2018}. Without sector-level evidence, investors, regulators, and audit committees are unable to determine whether audit and oversight elements are associated with reliable earnings.

\section{Research Gap}

Most research on audit quality and earnings management has been carried out in developed markets, where audit quality is usually measured by whether a company is audited by a member of a large international network. This measure does not suit the Pakistani setting well. \textcite{abid2018} report no significant difference in earnings management between companies audited by the Big Four and those audited by other firms. They attribute this to the small size of the audit market, low litigation risk and close economic ties between auditors and their clients.

In addition, most of the Pakistani studies reviewed in Chapter~2 combine companies from several industries in one sample, which makes their accrual estimates less comparable. Among these studies, only \textcite{alam2023}, who examine the cement sector, focus on a single industry. The Pakistani studies on audit tenure generally suggest a negative relationship with earnings management \parencite{kalabeke2019}, but they rely on multi-sector samples and examine periods before the 2019 Regulations came into force. None of the reviewed studies tests the direct link between audit fees and earnings management in Pakistan, and the findings on audit committee independence remain inconsistent.

As a result, none of the studies reviewed examines audit tenure, audit fees and audit committee independence together, within a single sector and under a single set of governance regulations. This study aims to address that gap.

\section{Research Questions}

\begin{enumerate}[label=\arabic*.]
  \item Is audit tenure associated with the level of earnings management in textile composite companies listed on the Pakistan Stock Exchange (PSX)?
  \item Are audit fees associated with the level of earnings management?
  \item Is greater audit committee independence associated with a lower level of earnings management?
  \item Do the findings remain consistent when signed discretionary accruals, lagged audit variables and the removal of outliers are applied?
\end{enumerate}

\section{Research Objectives}

\begin{enumerate}[label=\arabic*.]
  \item To examine whether audit tenure is related to earnings management.
  \item To examine whether audit fees are related to earnings management.
  \item To examine whether greater audit committee independence is related to lower earnings management.
  \item To assess the robustness of these findings.
\end{enumerate}

\section{Research Hypotheses}

\begin{description}[labelwidth=1cm,leftmargin=1.5cm,style=sameline]
  \item[H1:] There is a significant association between audit tenure and earnings management.
  \item[H2:] There is a significant association between audit fees and earnings management.
  \item[H3:] Greater audit committee independence is associated with lower earnings management.
\end{description}

\section{Scope of the Study}

This study is based on a purposive sample of the 20 largest ordinary share companies in the PSX Textile Composite sector. These companies are observed over six fiscal years, from 2020 to 2025, all of which fall under the 2019 Regulations, giving 120 company-year observations before data screening. Earnings management is measured using absolute discretionary accruals estimated through the Modified Jones Model, while the audit variables are audit tenure, audit fees and audit committee independence. Auditor size is not examined in this study. The sampling criteria are explained in Section~\ref{sec:sampling}, and the full list of companies is provided in Appendix~\ref{app:companies}.

\section{Significance of the Study}

The study is expected to benefit several groups. For investors and lenders, it offers a basis for judging how far reported profits can be relied upon. For audit committees and boards, it provides evidence on whether committee independence and audit arrangements are linked to lower earnings management. For regulators such as the SECP and PSX, as well as the audit profession, it shows how audit committee requirements and audit practices relate to the quality of financial reporting. Finally, for researchers, it adds single-sector evidence from an emerging market.

\needspace{10\baselineskip}% keep the short Limitations section on one page
\section{Limitations of the Study}
\label{sec:limitations}

Several limitations should be acknowledged. Audit tenure and audit fees act as proxies for audit quality rather than direct measures of it. Similarly, discretionary accruals capture earnings management only indirectly and are subject to estimation error. Annual reports may also disclose audit fees and auditor history inconsistently. The study identifies associations and cannot establish causation. Because it covers only one sector, its findings may not apply to other industries, and since the sample consists of the 20 largest companies in the sector rather than a random selection, the results may not reflect the position of smaller companies. Regulation 27 sets a minimum level of audit committee independence, which may limit variation in the independence ratio across companies. Auditor size is excluded, so differences between larger and smaller audit firms are not captured. Lastly, the study period is limited to six fiscal years.

\section{Structure of the Dissertation}

The dissertation is organized into five chapters. Chapter~1 introduces the background and sector context and sets out the problem, research gap, questions, objectives, hypotheses, scope, significance, and limitations. Chapter~2 reviews the theoretical and empirical literature and develops the hypotheses and conceptual framework. Chapter~3 explains the research design, sample, variables, and analytical techniques. Chapter~4 presents and discusses the results. Chapter~5 concludes the study by outlining its implications and recommendations and suggesting areas for future research.
